\documentclass[../main.tex]{subfiles}
\begin{document}

This chapter establishes the theoretical and empirical context for the thesis. It begins with the Big Five personality model and the NEO-PI-R facet schema that form the psychological foundation of the proposed system, then reviews the benchmark dataset used in this work. The subsequent sections trace the evolution of automatic personality recognition from linguistic feature engineering and deep learning through to modern large language model (LLM)-based approaches. The chapter concludes by identifying the limitations shared across existing methods, which directly motivate the design choices described in Chapter~\ref{chap:methodology}.

% =============================================================================
\section{The Big Five Personality Model}
\label{sec:bigfive}
% =============================================================================

Personality refers to a relatively stable set of individual differences in patterns of thought, emotion, and behaviour. Among the various personality frameworks proposed in psychology, the Big Five model --- also known as the OCEAN model --- has become the dominant empirical taxonomy in both academic research and applied settings~\cite{mccrae1992introduction,digman1990personality,tupes1992recurrent}. The model organises personality along five broad, statistically independent dimensions, summarised in Table~\ref{tab:big5}.

\begin{table}[ht]
  \centering
  \caption{The five dimensions of the Big Five (OCEAN) personality model.}
  \label{tab:big5}
  \begin{tabular}{p{2.8cm} p{9.5cm}}
    \toprule
    \textbf{Dimension} & \textbf{Description} \\
    \midrule
    Openness (O)        & Intellectual curiosity, aesthetic sensitivity, imagination, and receptiveness to new ideas and experiences. \\[4pt]
    Conscientiousness (C) & Tendency toward organisation, goal-directed discipline, reliability, and careful planning. \\[4pt]
    Extraversion (E)    & Sociability, assertiveness, positive emotionality, and a preference for external stimulation. Low scorers (introverts) prefer quieter, less stimulating environments. \\[4pt]
    Agreeableness (A)   & Cooperative, empathic, and trusting orientation toward others; associated with conflict avoidance and concern for others' welfare. \\[4pt]
    Neuroticism (N)     & Proneness to negative emotional states such as anxiety, depression, and irritability. Low scorers exhibit emotional stability. \\
    \bottomrule
  \end{tabular}
\end{table}

\subsection{The NEO-PI-R Facet Schema}

While the five broad dimensions provide a useful taxonomy, they are too coarse to capture within-trait variability that matters for personality analysis from text. Costa and McCrae addressed this with the Revised NEO Personality Inventory (NEO-PI-R)~\cite{costa2008revised}, which decomposes each of the five traits into six narrower \emph{facets}, giving thirty facets in total. The importance of this decomposition is best illustrated by example: two writers may share the same overall Extraversion score yet differ sharply --- one is gregarious and warmly sociable but passive, the other is assertive and energetic but prefers small groups. A single trait score discards this distinction; the six Extraversion facets preserve it. The thirty NEO-PI-R facets are listed in Table~3.1 of Chapter~\ref{chap:methodology}. This facet schema is the psychological representation used throughout the proposed system.

\subsection{Benchmark Datasets}

The standard evaluation corpus for automatic personality recognition from text is the \textbf{Essays dataset}~\cite{pennebaker1999linguistic}, consisting of 2,468 stream-of-consciousness essays written by university students, each labelled with binary (high/low) scores for all five OCEAN traits derived from a standardised self-report questionnaire. Because the essays are long, personal, and topically unconstrained, they are rich in psychologically relevant language but also in stylistic variation that can confound surface-level text models. This dataset is the primary benchmark used throughout this thesis.

% =============================================================================
\section{Traditional and Deep Learning Approaches}
\label{sec:traditional}
% =============================================================================

Research on automatic personality recognition predates the LLM era by more than two decades and has progressed through two broad methodological phases: hand-crafted linguistic features and end-to-end neural models.

\subsection{Linguistic Feature Engineering}

The earliest computational approaches were grounded in psycholinguistics. Pennebaker and colleagues developed the Linguistic Inquiry and Word Count (LIWC) tool~\cite{pennebaker2001linguistic}, which measures the frequency of words belonging to psychologically meaningful categories --- pronouns, positive emotion words, cognitive terms, and so on. The foundational finding of Pennebaker and King~\cite{pennebaker1999linguistic} was that writing style, independent of topical content, reliably correlates with personality: high-Extraversion writers use more social words, high-Conscientiousness writers use more deliberate and precise language. Mairesse and colleagues~\cite{mairesse2007using} systematically evaluated LIWC features in combination with classical classifiers such as Support Vector Machines (SVM) and established the first strong baselines on the Essays dataset.

The limitation of this approach is fundamental: hand-crafted word-category counts cannot capture contextual or compositional meaning. Two sentences sharing no vocabulary but conveying the same semantic content are treated as entirely unrelated, while two sentences sharing vocabulary but opposing in meaning are treated as similar. These failures become particularly acute for the interpersonal facets of Agreeableness and Extraversion, which tend to surface more through pragmatic implication than through explicit word choice.

\subsection{Deep Learning Models}

The shift to deep learning replaced manual feature design with representation learning directly from raw text. Majumder and colleagues~\cite{majumder2017deep} proposed splitting long essays into two halves, encoding each with a pair of CNNs over Word2Vec embeddings, and showed consistent improvements over LIWC-based models. Sun and colleagues~\cite{sun2018personality} combined Bidirectional LSTMs with CNN feature maps to capture both sequential and local structure in posts.

A notable step toward integrating psychological knowledge into neural models is PsyAttention~\cite{zhang2023psyattention}, published at EMNLP 2023. Although contemporaneous with early LLM-era work, PsyAttention remains a fine-tuned encoder model and is reviewed here as the strongest representative of the psycholinguistic feature + neural fusion paradigm. PsyAttention extracts 930 numerical psycholinguistic features from three text analysis tools --- SEANCE (sentiment and cognition), TAACO (cohesion), and TAALES (lexical sophistication) --- applies Pearson correlation filtering to reduce these to 108 non-redundant features on the Essays dataset, and encodes the resulting set with a Transformer-based attention encoder. This encoder is fused with a fine-tuned BERT encoder via a learned dynamic weighting layer. On the Essays dataset, PsyAttention achieves an average accuracy of 65.66\% across the five OCEAN traits, compared to 56.67\% for BERT alone. The result demonstrates that psycholinguistic features designed by domain experts still carry information that dense text representations miss, even when those representations come from large pre-trained models.

Despite this progress, deep learning models trained on the Essays or MyPersonality datasets face a common ceiling: the datasets are small by modern standards, training on them risks capturing dataset-specific stylistic patterns rather than personality-relevant content, and the learned representations do not generalise transparently across text genres.

% =============================================================================
\section{Large Language Model Approaches}
\label{sec:llm-methods}
% =============================================================================

The emergence of large language models such as GPT-3, GPT-4, and LLaMA has opened fundamentally different possibilities for personality recognition. Rather than training models from scratch on small labelled datasets, one can leverage the broad world knowledge and reasoning capabilities already encoded in LLMs. The field has explored several distinct strategies within this paradigm.

\subsection{LLMs as Text Augmenters}

A direct question when LLMs became available was whether they could simply outperform fine-tuned small models at personality classification. Hu and colleagues~\cite{hu2024llm} investigated this through a framework called TAE (Text Augmentation Enhanced). They found that when prompted directly to classify personality, ChatGPT performs comparably to --- but not better than --- fine-tuned BERT on MBTI datasets. The explanation is that LLMs focus on semantic content and may miss the subtler affective and stylistic cues that are diagnostically important for personality.

TAE therefore uses LLMs not as classifiers but as annotation assistants. For each post, an LLM is asked to generate three independent analyses --- semantic, sentiment, and linguistic --- which serve as positive pairs in a contrastive learning objective. The LLM also generates trait-specific explanations that enrich the label supervision signal. A smaller encoder trained with these augmented objectives outperforms fine-tuned BERT by 5.83\% on the Kaggle MBTI dataset and 6.53\% on the PANDORA dataset. The key lesson from TAE is that LLMs and small models are complementary: LLMs excel at synthesis and structured explanation, while small models exploit statistical patterns in labelled data.

\subsection{Psychological Questionnaires as Chain-of-Thought}

PsyCoT~\cite{yang2023psycot}, presented at EMNLP 2023, reframes personality classification as a structured reasoning task inspired by how psychologists actually assess personality. Rather than asking an LLM to predict a trait directly, PsyCoT uses the 44-item Big Five Inventory (BFI-44) as a chain-of-thought scaffold. In a multi-turn dialogue, the LLM is shown one questionnaire item per turn and asked to rate the author's text on that item on a 1--5 scale. After all 44 items have been rated, the accumulated ratings are aggregated by standard questionnaire scoring rules, and the LLM issues a final personality verdict informed by this history.

The multi-turn design is crucial: processing one item at a time keeps the model focused, and the historical ratings provide consistency constraints across related items. On the Essays dataset, PsyCoT improves average F1 by 4.23 points over standard prompting and surpasses or matches fine-tuned baselines on most traits. The method is sensitive to questionnaire granularity: replacing BFI-44 with the 10-item TIPI reduces performance substantially, confirming that the richer coverage of a 44-item inventory provides meaningfully more reasoning scaffolding. The main practical limitation is inference cost: 44 sequential LLM calls per sample is expensive at scale.

\subsection{Multi-Agent Comparative Assessment}

PADO~\cite{yeo2025pado} (Personality-induced multi-Agent framework for Detecting OCEAN), presented at COLING 2025, addresses a deeper challenge: personality is inherently relative and context-dependent, so a single agent's assessment may be unstable or biased. PADO introduces a three-stage pipeline. First, a \emph{personality induction} phase creates two reasoner agents with contrasting trait levels --- one prompted to embody high levels of the target trait and one low --- for each of the five OCEAN dimensions. Second, both agents analyse the author's text through three psycholinguistic lenses: emotional expression, cognitive style, and social aspects. Third, a \emph{judge agent} receives both explanations, compares them, and delivers a final verdict.

The comparative design exploits the relative nature of personality: rather than asking whether an author is extraverted in absolute terms, the judge agent asks which of two characterisations the text better supports. On the Essays dataset, PADO consistently outperforms zero-shot, one-shot, and Chain-of-Thought baselines for all tested LLMs (GPT-4o, GPT-3.5-turbo, Solar-10.7B-Instruct, LLaMA3-8B-Instruct). The gains are particularly pronounced for smaller models, where F1 scores increase by up to 0.37, suggesting that the structured comparison helps smaller models better leverage latent personality knowledge without additional fine-tuning.

\subsection{RAG-Based Personality Prediction}

PostToPersonality (P2P)~\cite{ma2025post}, presented at CIKM 2025, applies Retrieval-Augmented Generation to MBTI prediction from social media posts. The system operates in three steps. A locally fine-tuned LLM first reads a user's posts and produces a compact psychological assessment capturing key personality features. This assessment is then encoded into a vector and used to query a pre-built database of training-set posts, retrieving the $k$ most similar users together with their known MBTI labels. Finally, an online LLM receives the original posts, the local assessment, and the $k$ retrieved demonstrations in a single prompt and outputs the final MBTI prediction. To address the severe class imbalance inherent in the MBTI type distribution, P2P applies SMOTE oversampling during local LLM fine-tuning.

P2P outperforms 10 ML/DL baselines with an average F1 improvement of 20.17\% over the second-best method (XGBoost). The result demonstrates that RAG can effectively compensate for label scarcity and class imbalance in personality datasets. However, P2P retrieves using embeddings of the raw psychological assessment text, which means similarity is still measured over natural language rather than over a structured psychological representation. Two users with very different personality profiles but similar surface-level writing habits could still be retrieved as near neighbours.

\subsection{Ranking-Based Reinforcement Learning}

PerDet-R1~\cite{cao2026classification} takes a fundamentally different perspective on the problem formulation. Standard approaches treat the four binary MBTI dimensions independently, but the 16 MBTI types are not independent: INTJ and INTP share three dimensions, so conflating them is a much smaller error than conflating INTJ and ESFP. PerDet-R1 reframes personality prediction as a \emph{ranking} task --- the model must output the 16 types in order of compatibility with the user's posts --- and trains with Group Relative Policy Optimisation (GRPO) and an NDCG-based reward function that explicitly rewards correct ranking position. A first stage of supervised fine-tuning (SFT) with reasoning traces distilled from a larger teacher model bootstraps the policy before the RL stage begins. PerDet-R1 achieves state-of-the-art results on both the Kaggle and PANDORA benchmarks, demonstrating that modelling the interdimensional structure of personality types provides a meaningful advantage over treating each dimension as an independent binary classification.

Table~\ref{tab:related_summary} summarises the methods reviewed in this section.

\begin{table}[ht]
  \centering
  \caption{Summary of representative personality recognition approaches using LLMs.}
  \label{tab:related_summary}
  \begin{tabular}{p{2.4cm} p{3.0cm} p{3.5cm} p{4.2cm}}
    \toprule
    \textbf{Method} & \textbf{Core Idea} & \textbf{Strength} & \textbf{Limitation} \\
    \midrule
    TAE~\cite{hu2024llm}
      & LLM generates multi-perspective analyses as contrastive training data
      & Effective augmentation without labelled profiles; small model benefits from LLM knowledge
      & Relies on surface text similarity; no structured psychological representation \\[6pt]
    PsyCoT~\cite{yang2023psycot}
      & BFI-44 questionnaire items used as Chain-of-Thought steps
      & Structured reasoning; score aggregation provides consistency check
      & 44 LLM calls per trait; no labelled exemplars to calibrate reasoning \\[6pt]
    PADO~\cite{yeo2025pado}
      & Contrasting personality-induced agents plus a judge
      & Captures relative nature of personality; robust across model sizes
      & Multiple LLM calls; no training-set anchors; purely zero-shot \\[6pt]
    P2P~\cite{ma2025post}
      & Local LLM extracts features; RAG retrieves similar users by text
      & Handles class imbalance; demonstrated gains over strong baselines
      & Retrieval based on text similarity, not structured psychological profiles \\[6pt]
    PerDet-R1~\cite{cao2026classification}
      & Reframes prediction as type ranking; SFT + GRPO with NDCG reward
      & Models inter-type relationships; state-of-the-art on MBTI benchmarks
      & Targeted at MBTI ranking; requires RL training infrastructure \\
    \bottomrule
  \end{tabular}
\end{table}

% =============================================================================
\section{Limitations and Research Motivation}
\label{sec:motivation}
% =============================================================================

Reviewing the methods above reveals a limitation that none of them fully resolves: the difficulty of separating stable psychological content from surface writing style.

\paragraph{The surface style confound.} A model that directly embeds raw text --- whether through fine-tuned BERT or through LLM prompting --- learns from everything in that text simultaneously: topic, register, sentence length, emotional intensity, and personality-relevant content. Surface style is a confound because it correlates with personality in the training distribution but not necessarily across genres. An emotionally expressive introvert and an emotionally expressive extravert produce superficially similar text but carry opposite labels. Models that learn from raw text risk memorising this correlation rather than the underlying psychological signal. This problem is especially acute on the Essays dataset, where students' writing styles vary systematically by demographic and educational background in ways unrelated to personality.

\paragraph{Retrieval over raw text.} P2P demonstrates that RAG can be effective for personality prediction, but its retrieval operates over text-level features. A post about academic life written by a highly conscientious person and the same post written by a less conscientious person may produce similar retrieval vectors, because the topical content dominates the embedding. What is needed is retrieval over a psychological representation that is orthogonal to surface style.

\paragraph{Reasoning without labelled anchors.} PsyCoT and PADO both improve over standard prompting by introducing structured reasoning scaffolds, but neither provides the LLM with labelled training examples to calibrate against. Without such anchors, the LLM must rely entirely on its parametric knowledge of what personality traits look like in text --- knowledge that may be biased by the training distribution of the LLM itself.

These three observations jointly motivate the design of the proposed system. The central idea is to interpose an explicit psychological representation between raw text and the classifier: an LLM-generated thirty-facet profile in the NEO-PI-R vocabulary. Because this profile is expressed in a fixed psychological schema rather than in natural language, similarity between profiles measures psychological proximity rather than stylistic or topical overlap. Retrieval over profiles therefore returns training exemplars that share psychological structure with the test essay, not merely surface similarity. The retrieved exemplars --- rendered as full thirty-facet profiles with known labels --- then serve as calibration anchors that ground the final LLM reasoning step. The detailed design of this pipeline is presented in Chapter~\ref{chap:methodology}.

% =============================================================================
\section{Chapter Summary}
\label{sec:litreview-summary}
% =============================================================================

This chapter has traced the development of automatic personality recognition from psycholinguistic feature engineering through deep learning to modern LLM-based methods. The Big Five model and NEO-PI-R facet schema provide a principled psychological foundation; the Essays dataset establishes the empirical benchmark. Traditional methods revealed the importance of psycholinguistic signals; deep learning models showed that neural representations can surpass hand-crafted features but remain limited by small labelled datasets. LLM-based approaches --- TAE, PsyCoT, PADO, P2P, and PerDet-R1 --- each address a different aspect of the problem, yet all share the challenge of disentangling surface style from psychological content. This shared limitation motivates the proposed three-component pipeline: a psychological profiler, a profile-based hybrid retriever, and a retrieval-augmented reasoner.

\end{document}
